\documentclass[10pt]{book}
\usepackage[T1]{fontenc}
\usepackage[utf8]{inputenc}
\usepackage[french]{babel}
\usepackage{amsmath,amssymb}
\usepackage[colorlinks=true,hypertexnames=false,unicode=true,pdfencoding=auto]{hyperref}
\newcommand{\sh}{\mathcal}
\begin{document}
\hypertarget{ega4.c04.subsec.16-1}{}\hypertarget{ega4.c04.subsec.16-4}{}
\begingroup
\small
\newcommand{\EGAIVcontentsline}[3]{%
  \par\begingroup
  \rightskip=2.5em\parfillskip=-\rightskip
  \noindent #1\if\relax\detokenize{#1}\relax\else.\fi\quad #2\nobreak\hspace{0.5em}%
  \dotfill\nobreak\hbox to 2.5em{\hss #3}\par
  \endgroup}
\hfill\textsc{Pages}\par

\EGAIVcontentsline{}{\textsc{Chapitre IV.} --- \textbf{Étude locale des schémas et des morphismes de schémas} \textit{(fin)}}{5}

\EGAIVcontentsline{\S16}{Invariants différentiels. Morphismes différentiellement lisses}{5}
\EGAIVcontentsline{{\texorpdfstring{\protect\hyperlink{ega4.c04.subsec.16-1}{16.1}}{16.1}}}{Invariants normaux d'une immersion}{5}
\EGAIVcontentsline{{\texorpdfstring{\protect\hyperlink{ega4.c04.subsec.16-4}{16 4}}{16.4}}}{Propriétés fonctorielles des invariants différentiels}{16}
\EGAIVcontentsline{{\texorpdfstring{\protect\hyperlink{ega4.c04.subsec.16-7}{16.7}}{16.7}}}{Les $\sh{P}^{n}_{X/S}(\sh{F})$}{36}
\EGAIVcontentsline{{\texorpdfstring{\protect\hyperlink{ega4.c04.subsec.21-14}{21.14}}{21.14}}}{Le théorème de Ramanujam-Samuel}{323}
\EGAIVcontentsline{}{\textsc{Bibliographie}}{333}
\EGAIVcontentsline{}{\textsc{Index terminologique}}{336}

\begin{center}\rule{25mm}{0.4pt}\end{center}

\endgroup

\clearpage
% La page imprimée 344 est blanche.
\marginpar{\scriptsize\textbf{IV}~|~344}\ignorespaces
\thispagestyle{empty}
\null
\clearpage
x
\end{document}
